\section{Limitations, and where each algorithm is worth nothing}
\label{sec:limits}

Collected in one place so that nothing in this document has to be read charitably.

\subsection{Algorithm~\ref{alg:gfp} (liveness)}

\begin{description}
  \item[Zeno-ness is inherited, not solved.] The cycle witness is not divergence-checked. A
        purpose-built discriminator (\S\ref{sec:soundness:zeno}) shows \verifyta{}, the reference
        engine and Algorithm~\ref{alg:gfp} \emph{all} report the Zeno run as a violation, so this is
        not a defect relative to the oracle --- but it is a limit of all three, and a
        divergence-aware semantics would need the standard fresh-clock construction added to each.
  \item[Partial deadlocks.] $\mathrm{stop}(v)$ uses the engine's exact partial-deadlock routine as a
        per-node flag, which over-approximates the stuck \emph{set}. The reference engine has its own
        documented incompleteness here. Both sides of every comparison share it, so the differential
        results cannot detect it.
  \item[Federation width is not bounded by any theorem.] It was $1.00$--$1.35$ everywhere measured.
        \texttt{grc4} initially exceeded a cap of $256$; raising the cap to $4096$ made it converge
        at width $1.00$, so the cap --- not the algorithm --- was the limit. But there is still no
        proof that width stays small, only an observation over seven families.
  \item[Worthless on violated queries.] The existing engine answers those within a handful of nodes
        because counterexamples sit near the initial state. Algorithm~\ref{alg:gfp} must complete
        the forward phase first and is strictly worse there. It pays only on \emph{satisfied}
        liveness queries.
  \item[Says nothing about \texttt{A[]}.] The reachability floor stands untouched, and
        \S\ref{sec:deferral} closes the \texttt{explored} axis for uniform policies.
\end{description}

\subsection{The reachability side (\S\ref{sec:deferral})}

There is no algorithm here to limit; deferral was implemented, measured, and found to save exactly
zero. What remains are limits on the \emph{negative}:

\begin{description}
  \item[The negative is about UNIFORM deferral.] Proposition~\ref{prop:uniform} closes uniform
        policies under any schedule. \textbf{Selective} deferral is untouched and remains the only
        open door --- it needs a per-state predictor, and we showed only that no
        \emph{location-vector} predictor exists on \texttt{fischer5}. A predictor built from
        something else (edge history, arrival order, a cheap zone summary) is not excluded.
  \item[One interleaving.] The barrier-free frontier was serialised as a global FIFO. That is one
        asynchronous interleaving, not all of them; a genuinely concurrent implementation could in
        principle behave differently, though Proposition~\ref{prop:uniform} says it should not.
  \item[The probe is not the engine.] No convex-union merge, so absolute counts differ, and the
        \texttt{phil} conclusion (zero cross-level retirements) cannot be transferred to the engine,
        which does merge and where the prior study measured a cascade on \texttt{phil12}.
  \item[``Independent'' is about the shortest path.] The BFS tree stores one parent per node, so the
        $100\%$ figure means an independent discovery path \emph{exists}, not that every path is.
\end{description}

\subsection{Threats to the evaluation itself}

\begin{itemize}
  \item The oracle is the shipping engine, not a proof. Where prototype and engine disagreed,
        \verifyta{} broke the tie --- but two implementations sharing a semantics can share a bug,
        and \texttt{grc4}'s partial-deadlock queries are a known instance where the engine and
        \verifyta{} already disagree with each other.
  \item Seven model families, small instances. The instances are small precisely because the
        \emph{reference} side does not scale, which is the phenomenon under study; that makes the
        comparison honest but the sample narrow.
  \item Two query shapes (\texttt{p --> q}, 164/164; \texttt{E[]}, 51/51). \texttt{A<>} follows by
        negation but was not run separately. \textbf{Nested} formulas are untested, and the
        least-fixpoint half of the logic (\texttt{E<>}, \texttt{A[]}) is untouched by this work ---
        it is already at the floor.
\end{itemize}

\subsection{The single most likely way this is wrong}

If federation width blows up on a family not yet tried, Algorithm~\ref{alg:gfp} degrades towards
the status quo rather than failing outright --- that is the benign failure. The malign one is the
partial-deadlock flag: it is used as a per-node flag and therefore over-approximates the stuck set,
so on a model where it fires spuriously the algorithm will report a violation that is not there.
Every other known gap in this document is in the safe direction (missing a violation, not inventing
one). The cheap experiment is to widen the differential suite beyond the current seven families.
